\section*{第 \ref{ch:b1:heat-engines} 章——热机与机器}

\begin{solution}{exo:b1:heat-engines:1}
$\eta_C = 1 - 300/800 = 0.625$。取 $0.40$ 时：$\dot Q_h = 1.0/0.40 = \qty{2.5}{GW}$
从燃料中吸热，$\dot Q_c = 2.5 - 1.0 = \qty{1.5}{GW}$ 排入河中——
排掉的热量比售出的电能还多。
\end{solution}

\begin{solution}{exo:b1:heat-engines:2}
$e_f \leq 275/(298 - 275) = 12$。取 $e_f = 3$ 时：$P = 100/3 = \qty{33}{W}$
（另有 \qty{133}{W} 进入厨房）。
\end{solution}

\begin{solution}{exo:b1:heat-engines:3}
$r = 9$：$9^{0.4} = 2.41$，$\eta = 0.585$；$r = 11$：$11^{0.4} = 2.61$，
$\eta = 0.62$。当 $r = 20$ 时，汽油－空气混合物被压缩到约
\qty{900}{K}，会在火花塞点火之前自燃（爆震），损坏
发动机；只有仅压缩空气的柴油机才能达到这一压缩比。
\end{solution}

\begin{solution}{exo:b1:heat-engines:4}
$e_p \leq 300/20 = 15$：理想功率 $5/15 = \qty{0.33}{kW}$；实际功率 $5/3.5 =
\qty{1.43}{kW}$；电阻器则为 \qty{5}{kW}。
\end{solution}

\begin{solution}{exo:b1:heat-engines:5}
绝热线：$V_C = V_B(T_h/T_c)^{1/(\gamma-1)} = 3.0 \times (5/3)^{2.5} =
\qty{10.8}{L}$，$V_D = 1.0 \times 3.59 = \qty{3.59}{L}$（故 $V_C/V_D = 3 =
V_B/V_A$）。$Q_h = RT_h\ln3 = 8.314 \times 500 \times 1.099 = \qty{4.57}{kJ}$，
$Q_c = -RT_c\ln3 = \qty{-2.74}{kJ}$，$W = -(4.57 - 2.74) = \qty{-1.83}{kJ}$，
$\eta = 1.83/4.57 = 0.40 = 1 - 300/500$。
\end{solution}

\begin{solution}{exo:b1:heat-engines:6}
先看第一种说法：能量 $200 = 1000 - 800$ 平衡；克劳修斯不等式 $1000/400 - 800/300
= 2.5 - 2.67 = \qty{-0.17}{J/K}$，故 $S_{\mathrm{created}} = \qty{0.17}{J/K} \geq 0$，
可能成立（$\eta = 0.20$）。再看第二种：$1000/400 - 700/300 = +\qty{0.17}{J/K} > 0$，
意味着熵被消灭了——不可能。最大效率：$\eta_C = 1 -
300/400 = 0.25$，即 \qty{250}{J}。
\end{solution}

\begin{solution}{exo:b1:heat-engines:7}
$r^{\gamma-1} = 18^{0.4} = 3.18$；$\rho^\gamma = 2^{1.4} = 2.64$：$\eta = 1 -
1.64/(1.4 \times 3.18 \times 1) = 0.63$。奥托循环在 $r = 18$ 时：$1 - 1/3.18 = 0.69$；
在 $r = 9$ 时：$0.585$。柴油机压缩的是纯空气：压缩过程中不存在燃料，无法提前点火，
燃料喷入 \qty{900}{K} 的空气后即刻燃烧，
因此不存在爆震极限。
\end{solution}

\begin{solution}{exo:b1:heat-engines:8}
$Q_h = RT_h\ln2 = 8.314 \times 600 \times 0.693 = \qty{3.46}{kJ}$，$Q_c = -8.314
\times 300 \times 0.693 = \qty{-1.73}{kJ}$：每个循环 $W = \qty{-1.73}{kJ}$。有
回热器时 $\eta = 1.73/3.46 = 0.50 = 1 - 300/600$。无回热器时，热源还须
提供 $C_{V,m}(T_h - T_c) = 2.5 \times 8.314 \times 300 =
\qty{6.24}{kJ}$：$\eta = 1.73/(3.46 + 6.24) = 0.18$。
\end{solution}

\begin{solution}{exo:b1:heat-engines:9}
两个物块的熵变：$\Delta S = C\ln(T_f/T_1) + C\ln(T_f/T_2) =
C\ln\bigl(T_f^2/T_1T_2\bigr) = S_{\mathrm{created}} \geq 0$（热机自身的
熵在每个循环后都回到原值），故 $T_f \geq \sqrt{T_1T_2}$，且
$W = C(T_1 + T_2 - 2T_f)$（第一定律）在 $T_f$ 最小时最大：
可逆情形下 $T_f = \sqrt{400 \times 300} = \qty{346.4}{K}$，
$W_{\max} = 4000 \times (700 - 692.8) = \qty{29}{kJ}$。直接接触时：$T_f =
\qty{350}{K}$，不做功，$S_{\mathrm{created}} = 4000\ln(350^2/120000) =
\qty{82}{J/K}$。
\end{solution}

\begin{solution}{exo:b1:heat-engines:10}
(a) $P = \dot Q_h - |\dot Q_c| = K(T_h - x) - K(y - T_c)$；可逆的核心
不产生净熵交换：$K(T_h - x)/x = K(y - T_c)/y$。(b) 交叉相乘：
$y(T_h - x) = x(y - T_c)$，故 $y(T_h - 2x) = -xT_c$，$y = xT_c/(2x - T_h)$。
于是 $P = K(T_h - x)(1 - y/x) = K(T_h - x)(2x - T_h - T_c)/(2x - T_h)$；令
$u = 2x - T_h$，则 $T_h - x = (T_h - u)/2$，且 $P = K(T_h - u)(u - T_c)/(2u) =
\tfrac{K}{2}[(T_h + T_c) - u - T_hT_c/u]$。(c) $\dd P/\dd u = \tfrac K2(-1 +
T_hT_c/u^2) = 0$，解得 $u = \sqrt{T_hT_c}$，为极大值；此时 $x = (T_h +
\sqrt{T_hT_c})/2$，$y = xT_c/u = (\sqrt{T_hT_c} + T_c)/2$，且 $\eta^\ast = 1 - y/x
= 1 - \sqrt{T_c}(\sqrt{T_h} + \sqrt{T_c})/[\sqrt{T_h}(\sqrt{T_h} + \sqrt{T_c})] = 1 -
\sqrt{T_c/T_h}$；$P_{\max} = \tfrac K2(\sqrt{T_h} - \sqrt{T_c})^2$。(d)
$\eta^\ast = 1 - \sqrt{0.375} = 0.39$，而 $\eta_C = 0.625$：实际值 $0.40$
（\cref{exo:b1:heat-engines:1}）所示的热机是为功率而设计的，
而非为效率。
\end{solution}

\begin{solution}{exo:b1:heat-engines:11}
循环过程：$Q_h + Q_c + Q_0 = 0$，且 $Q_h/T_h + Q_c/T_c + Q_0/T_0 \leq 0$。
消去 $Q_0 = -(Q_h + Q_c)$：$Q_h(1/T_h - 1/T_0) + Q_c(1/T_c - 1/T_0) \leq 0$，
即 $Q_c(T_0 - T_c)/(T_cT_0) \leq Q_h(T_h - T_0)/(T_hT_0)$，从而
$Q_c/Q_h \leq (1 - T_0/T_h)\,T_c/(T_0 - T_c)$。第一个因子是
燃烧器与室温之间卡诺热机的效率，第二个因子是
舱室与室温之间卡诺制冷机的性能系数：该上界
即两者理想耦合所能得到的结果。数值：$0.25 \times 11 = 2.75$
（实际吸收式制冷机约为 $0.6$）。
\end{solution}

\begin{solution}{exo:b1:heat-engines:12}
(a) \qty{0.40}{J}。(b) $0.40 \times 3.5 = \qty{1.4}{J}$。(c) \qty{0.90}{J}。
(d) 热机 $\eta = 1 - 293/800 = 0.634$，热泵 $e_p = 293/20 = 14.65$：
$0.634 \times 14.65 = \qty{9.3}{J}$（热机自身的 \qty{0.37}{J}
在 \qty{293}{K} 排出，也可用来供暖：共计 \qty{9.7}{J}）。多于一焦耳，
是因为其余部分取自室外空气：能量守恒，
而热力学第二定律只要求熵不被消灭，
可逆的串联系统恰好满足这一要求。
\end{solution}

\begin{solution}{pb:b1:heat-engines:1}
\textbf{1.} $P = 250 \times 25 = \qty{6.25}{kW}$。

\textbf{2.} 平均损耗 $250 \times 13 = \qty{3.25}{kW}$，对应 $200 \times 86400 =
\qty{1.73e7}{s}$：$E = \qty{5.6e10}{J} = \qty{15600}{kW h}$。

\textbf{3.} $15600 \times 0.25 = 3900$~\texteuro.

\textbf{4.} 燃气 $15600/0.90 = \qty{17300}{kW h}$：$1730$~\texteuro。

\textbf{5.} 电阻器：$15600/0.40 = \qty{39000}{kW h}$ 的燃料；锅炉
\qty{17300}{kW h}：少 $2.25$ 倍。电能是功，是能做任何事的
形式；把它在电阻器中变为 \qty{20}{\degreeCelsius} 的热量，
就舍弃了电厂中已损失的 $60\%$ 燃料，以及
泵送热量的可能性。

\textbf{6.} 一个循环中 $W + Q_h + Q_c = 0$，且 $Q_h/T_{\mathrm{in}} +
Q_c/T_{\mathrm{out}} \leq 0$，其中 $Q_h < 0$、$Q_c > 0$：$Q_c \leq
|Q_h|T_{\mathrm{out}}/T_{\mathrm{in}}$，于是 $W = |Q_h| - Q_c \geq |Q_h|(1 -
T_{\mathrm{out}}/T_{\mathrm{in}})$，且 $e_p = |Q_h|/W \leq T_{\mathrm{in}}/(T_{\mathrm{in}} -
T_{\mathrm{out}})$。

\textbf{7.} $293/25 = 11.7$; $6.25/11.7 = \qty{0.53}{kW}$.

\textbf{8.} $293/13 = 22.5$；$15600/22.5 = \qty{690}{kW h}$（约
$170$~\texteuro）。

\textbf{9.} 每提升一焦耳所做的功为 $(T_{\mathrm{in}} - T_{\mathrm{out}})/
T_{\mathrm{in}}$，即提升的“高度”；室外越冷，
每焦耳须泵得越高——所需的焦耳也越多，因为
损耗也正比于 $T_{\mathrm{in}} - T_{\mathrm{out}}$：
电功率随温差的平方增长，
恰在供暖最要紧的时候最糟。

\textbf{10.} 取自被冷却的室外空气：$Q_c = 6.25 - 0.53 =
\qty{5.7}{kW}$。热力学第二定律禁止热量\emph{自发地}
从冷处流向热处；这里由功来支付“提升”。核对：从空气中取走的熵
$5.72/268 = \qty{21.3}{W/K}$，交给房屋的熵 $6.25/293 =
\qty{21.3}{W/K}$——无熵被消灭，也无熵被产生：理想的热泵。

\textbf{11.} 压缩机（接收 $W$）、冷凝器（放出 $|Q_h|$），
膨胀阀（两者皆不），蒸发器（接收 $Q_c$）。热量只会
沿温度梯度向下流动：制冷剂须比空气更冷，才能从空气取热，
并须比供暖水更热，才能向它放热。

\textbf{12.} $T_{\mathrm{ev}} = 268 - 8 = \qty{260}{K}$; $e_{p,C} = 313/(313 - 260)
= 5.9$.

\textbf{13.} $e_p = 0.55 \times 5.9 = 3.25$; $W = 6.25/3.25 = \qty{1.92}{kW}$;
$Q_c = 6.25 - 1.92 = \qty{4.33}{kW}$，取自空气。

\textbf{14.} $\dot m = 4330/200000 = \qty{22}{g/s}$，每分钟 \qty{1.3}{kg}。

\textbf{15.} $T_{\mathrm{ev}} = \qty{272}{K}$, $313/41 = 7.6$, $e_p = 4.2$;
$W = 3.25/4.2 = \qty{0.77}{kW}$.

\textbf{16.} $e_p = 0.55 \times 313/(321 - T_{\mathrm{out}}) = 172/(321 -
T_{\mathrm{out}})$（$T_{\mathrm{out}}$ 以开尔文计）。季节：$15600/4.2 =
\qty{3700}{kW h}$，$930$~\texteuro——对比电阻器的 $3900$ 与
燃气的 $1730$。

\textbf{17.} $3700/0.40 = \qty{9300}{kW h}$ 的燃料：比锅炉少 $46\%$，
比电阻器少四倍，即便经过效率为 $40\%$ 的
发电厂。

\textbf{18.} 需求 $250 \times 35 = \qty{8.75}{kW}$；$T_{\mathrm{ev}} = \qty{250}{K}$，
$313/63 = 5.0$，$e_p = 2.7$，$W = \qty{3.2}{kW}$（若能如此）。容量
\qty{3.8}{kW}：备用 $8.75 - 3.8 \approx \qty{5}{kW}$ 由电阻器提供。
按最冷一小时选型的热泵会过大——而且全冬都在
低效地循环运行；应按设计温度选型，
让廉价的电阻器承担罕见的极端情况。

\textbf{19.} $338/(338 - 260) = 4.33$, $e_p = 2.4$, $W = 6.25/2.4 =
\qty{2.6}{kW}$：比地板温度为
\qty{35}{\degreeCelsius} 时多耗 $36\%$ 的电能。低温散热末端占
安装费用的一半。

\textbf{20.} 一个循环中的熵：$Q_h/T_{\mathrm{in}} + Q_c/T_{\mathrm{out}} +
S_{\mathrm{created}} = 0$ 得 $Q_c = T_{\mathrm{out}}|Q_h|/T_{\mathrm{in}} -
T_{\mathrm{out}}S_{\mathrm{created}}$；代入 $W = |Q_h| - Q_c$。

\textbf{21.} 冷凝器：$6250\,(1/293 - 1/313) = \qty{1.36}{W/K}$；蒸发器：
$4330\,(1/260 - 1/268) = \qty{0.50}{W/K}$.

\textbf{22.} $268 \times 1.86 = \qty{0.50}{kW}$：恰为
\qty{1.92}{kW} 的四分之一。

\textbf{23.} $(1.92 - 0.53)/268 = \qty{5.2}{W/K}$；换热器贡献
$36\%$；其余来自压缩机（摩擦、非等熵
压缩、电动机损耗）、阀门中的节流（
自由膨胀）、压降、风扇。

\textbf{24.} $T_{\mathrm{cd}} = \qty{303}{K}$，$T_{\mathrm{ev}} = \qty{264}{K}$：
冷凝器 $6250\,(1/293 - 1/303) = \qty{0.70}{W/K}$，蒸发器
$4330\,(1/264 - 1/268) = \qty{0.25}{W/K}$：共 \qty{0.95}{W/K}，而非
$1.86$，即节省 $268 \times 0.91 \approx \qty{0.25}{kW}$（$13\%$；
整台机器以更优卡诺值的 $0.55$ 计，将获得
更多收益）。温差减半，意味着换热面积加倍，换热器更大，
更贵，风扇功耗也更多。

\textbf{25.} $15600/3700 = 4.2$ 千瓦时热量，每千瓦时
电能；$930$~\texteuro\ 对比燃气的 $1730$；碳排放：$3700
\times 0.060 = \qty{220}{kg}$ 对比 $17300 \times 0.20 = \qty{3500}{kg}$——
少十五倍。
\end{solution}
